\\documentclass[11pt]{article}

\\usepackage[margin=1in]{geometry}
\\usepackage{amsmath,amssymb}
\\usepackage{booktabs}

\\title{Problem Statement: Block-Based Active Noise Control Simulator}
\\author{}
\\date{}

\\begin{document}
\\maketitle

\\section*{Objective}
This simulator models a feedforward active noise control (ANC) system driven by a monaural WAV file. The purpose is to generate two microphone signals, an outside microphone and an in-ear microphone, under a simple acoustic plant model, while allowing a user-supplied control law to produce anti-noise in real time on a block-by-block basis.

\\section*{Signal Model}
All signals are processed in blocks of
\\[
N = 256
\\]
samples at a sample rate of
\\[
f_s = 48{,}000 \, \text{Hz}.
\\]
Each acoustic path is represented by a finite impulse response (FIR) filter of length
\\[
M = 1024.
\\]

Let $k$ denote the block index. The simulator uses the following signals:
\\begin{align*}
x_k & : \text{input audio block from the WAV source}, \\
v_k & : \text{synthetic additive disturbance block}, \\
n_k & = x_k + v_k : \text{total ambient disturbance}, \\
u_k & : \text{controller output block}, \\
u_k^{\mathrm{spk}} & : \text{speaker-shaped control block}, \\
y_k^{H} & : \text{disturbance at the outside microphone}, \\
y_k^{P} & : \text{disturbance at the in-ear microphone}, \\
y_k^{C} & : \text{control leakage to the outside microphone}, \\
y_k^{S} & : \text{control contribution at the in-ear microphone}.
\\end{align*}

The simulator contains four fixed FIR paths and one time-varying FIR path:
\\begin{align*}
H & : \text{noise} \rightarrow \text{outside microphone}, \\
P & : \text{noise} \rightarrow \text{in-ear microphone}, \\
C & : \text{speaker} \rightarrow \text{outside microphone}, \\
G & : \text{speaker response}, \\
S_k & : \text{speaker} \rightarrow \text{in-ear microphone}.
\\end{align*}

\\section*{Plant Equations}
At every block, the plant is propagated according to
\\begin{align}
u_k^{\mathrm{spk}} &= G * u_k, \\
y_k^{C} &= C * u_k^{\mathrm{spk}}, \\
y_k^{S} &= S_k * u_k^{\mathrm{spk}}, \\
y_k^{H} &= H * n_k, \\
y_k^{P} &= P * n_k,
\\end{align}
where $*$ denotes linear convolution implemented blockwise.

The microphone signals presented by the simulator are then
\\begin{align}
m_k^{\mathrm{out}} &= y_k^{H} + y_k^{C}, \\
m_k^{\mathrm{in}} &= y_k^{P} + y_k^{S}.
\\end{align}

The outside microphone therefore observes the disturbance plus speaker leakage, while the in-ear microphone observes the disturbance plus the anti-noise delivered through the secondary path.

\\section*{Latency Model}
The controller does not act instantaneously. Instead, the simulator applies a configurable delay of $L$ blocks between controller computation and plant actuation. In the current configuration,
\\[
L = 1.
\\]
Thus, a control block computed from microphone block $k$ is injected into the plant after one block of delay.

\\section*{Secondary-Path Variation}
The in-ear secondary path is allowed to vary over time in order to emulate mild plant drift. Let $S_0$ denote the nominal secondary path. For each block, the simulator generates a filtered random perturbation $w_k$ and forms
\\[
\widetilde{S}_k = S_0 + \gamma w_k,
\\]
where $\gamma$ is the parameter \texttt{noise\_gain}. The perturbed path is then renormalized to preserve the nominal path norm,
\\[
S_k = \frac{\lVert S_0 \rVert_2}{\lVert \widetilde{S}_k \rVert_2} \, \widetilde{S}_k,
\\]
and each coefficient is clipped to the interval $[-1,1]$.

\\section*{Synthetic Noise Model}
In addition to the WAV input, the simulator adds a colored random disturbance. A random cutoff-frequency trajectory is generated, low-pass filtered, and used to parameterize samplewise Gaussian draws. The resulting block is mean-removed and peak-normalized to a prescribed scale before being added to the WAV input. This produces a disturbance that is random, slowly varying, and spectrally colored.

\\section*{Controller Interface}
At block $k$, the controller receives the microphone pair
\\[
\left(m_k^{\mathrm{out}},\; m_k^{\mathrm{in}}\right)
\\]
and must produce a control block $u_k$. The reference implementation currently uses the simple baseline law
\\[
u_k = -m_k^{\mathrm{in}},
\\]
which is intended only as a placeholder controller.

\\section*{Problem Statement}
Given the disturbance input, the acoustic paths $H$, $P$, $C$, $G$, and the time-varying secondary path $S_k$, design a block-based controller that computes $u_k$ from the available microphone measurements so that the in-ear signal energy
\\[
\sum_k \lVert m_k^{\mathrm{in}} \rVert_2^2
\\]
is minimized while remaining causal and respecting the one-block actuation delay.

The simulator should therefore be understood as a discrete-time experimental platform for testing ANC algorithms under convolutional acoustic propagation, additive colored noise, mild plant variation, and explicit control latency.

\\end{document}
